\section{Introduzione della previsione}
\label{sec:model}
\subsection{Che cosa deve imparare il modello}
Prima di usare il modello su una nuova mappa, lo alleno su esempi di
percorsi parziali. Questa preparazione è detta fase \emph{offline}.
La fase successiva, in cui uso il modello già allenato per cercare un giro
su una mappa, è detta \emph{ricerca online}. Qui ``online'' indica la fase
di utilizzo del modello, non una connessione a Internet.

Il modello deve prevedere quanta strada manca per visitare le consegne
residue e tornare al deposito. Chiamo questa previsione $\widehat h(s)$.
La distanza $g$ già percorsa è nota e viene tenuta separata.
Il valore ideale sarebbe il costo residuo ottimo, indicato con $h^*(s)$.
Sulle mappe grandi, però, gli esempi di allenamento usano completamenti
approssimati: la previsione cerca di riprodurre quei costi di riferimento.

\subsection{Le informazioni usate per prevedere}
Per descrivere ogni stato al modello scelgo quattro numeri, chiamati
\emph{feature}. Riassumono quanto resta da fare e la geometria delle prossime
visite. Come nella Sezione~1, $v$ è il luogo corrente e $N$ l'insieme delle
consegne residue. Chiamo $m$ il loro numero, scritto anche $m=|N|$.

\begin{center}\small
\begin{tabular}{lll}
\toprule
Simbolo & Informazione & Calcolo\\
\midrule
$x_1$ & Numero di consegne residue & $m$\\
$x_2$ & Distanza dal luogo corrente al deposito & $d(v,0)$\\
$x_3$ & Distanza dalla consegna residua più vicina & $\min_{a\in N}d(v,a)$\\
$x_4$ & Media delle distanze dalle consegne residue & $\frac{1}{m}\sum_{a\in N}d(v,a)$\\
\bottomrule
\end{tabular}
\end{center}
Se non restano consegne, minimo e media sono posti a zero.

Durante la ricerca, quando tutte le consegne sono state visitate,
\textbf{non serve prevedere la strada mancante}: l'unica tappa è il ritorno
al deposito, e la sua distanza si legge direttamente dalla mappa.

\subsection{Come viene calcolata la previsione}
Uso una regressione lineare: il modello moltiplica ciascuna feature per
un peso e somma i quattro contributi a un termine costante. Indico con
$w_1,\ldots,w_4$ i pesi e con $b$ il termine costante, detto intercetta.
La previsione è quindi
\[
 \widehat h(s)=b+\sum_{j=1}^{4}w_jx_j(s).
\]
Nell'allenamento conosco, per ogni esempio, un costo rimanente di riferimento,
chiamato \emph{target} e indicato con $y(s)$. I pesi sono scelti per rendere
piccole le differenze fra target e previsioni. Chiamando $D_{train}$
l'insieme degli esempi di allenamento, il criterio è
\[
 \min_{b,w}\sum_{s\in D_{train}}\bigl(\widehat h(s)-y(s)\bigr)^2.
\]
Questo criterio determina i pesi e l'intercetta minimizzando l'errore
quadratico sui dati di allenamento. Le feature vengono usate nelle loro
unità originali, senza normalizzarle.

\subsection{Parametri appresi e limiti}
Per il modello da cento consegne, l'intercetta e i quattro pesi risultano,
arrotondati per la lettura:
\[
\begin{split}
 b&=163{,}645698,\\
 (w_1,w_2,w_3,w_4)&=(6{,}027693;\ 0{,}017960;\ -1{,}597434;\ 1{,}667937).
\end{split}
\]
Il peso negativo di $x_3$ è il risultato dell'allenamento con informazioni
fra loro correlate. Non significa che allontanare una consegna renda
necessariamente più breve il giro.

Queste quattro informazioni non descrivono tutta la disposizione dei punti:
stati diversi possono avere le stesse feature e costi residui differenti.
Anche il termine costante può pesare molto quando restano poche consegne.
Il modello può quindi sovrastimare o sottostimare il costo rimanente:
la sua previsione aiuta a scegliere l'ordine di esplorazione, ma non viene
usata per scartare definitivamente un ramo.
